\documentclass[12pt]{article}
\usepackage[a4paper, margin=2cm]{geometry}
\usepackage{amsmath}
\usepackage{array}
\usepackage{tikz}
\usepackage{tikz-feynman}
\tikzfeynmanset{compat=1.1.0}

\setlength{\parindent}{0pt}

\begin{document}

\section*{Weak Interactions -- Initial Stage}

\textbf{Worked example:} A strange quark changes into an up quark by emitting a \(W^-\) boson:
\[
s \rightarrow u + W^-
\]

\vspace{0.5cm}

% Two-column layout
\begin{tabular}{l l}

% LEFT COLUMN
\begin{minipage}[t]{0.46\textwidth}
\vspace{0pt}
\textbf{Feynman diagram}

\vspace{0.6cm}

% --- TEMPLATE (labels changed only) ---
\begin{tikzpicture}
\begin{feynman}
  % Vertices
  \vertex (l1) at (0,0) {\(s\)};
  \vertex [coordinate] (a) at (1,2) {};
  \vertex (l2) at (0,4) {\(u\)};
  \vertex (b) at (4,2) {};

  \diagram*{
    (l1) -- [fermion] (a) -- [fermion] (l2),
    (a) -- [boson, edge label=\(W^-\), momentum'={}] (b),
  };
\end{feynman}
\end{tikzpicture}
% --- END TEMPLATE ---
\end{minipage}
&
% RIGHT COLUMN
\begin{minipage}[t]{0.50\textwidth}
\vspace{0pt}

\renewcommand{\arraystretch}{1.4}
\begin{tabular}{@{}l c c c c c@{}}
 & \(s\) & \(\Rightarrow\) & \(u\) & \(+\) & \(W^-\) \\
Charge & $-\tfrac{1}{3}$ & $=$ & $+\tfrac{2}{3}$ & $+$ & $-1$ \\
Lepton no. & $0$ & $=$ & $0$ & $+$ & $0$ \\
Baryon no. & $\tfrac{1}{3}$ & $=$ & $\tfrac{1}{3}$ & $+$ & $0$ \\
\end{tabular}

\end{minipage}

\end{tabular}

\vspace{0.8cm}

\textbf{Conclusion:} Charge, lepton number, and baryon number are conserved, so this flavour change is allowed in the weak interaction.

\vspace{0.9cm}

\section*{Practice: Weak Interactions -- Flavour Change}

%%%%%%%%%%%%%%%%%%%%%%%%%%%%%%%%%%%%%%%%%%%%
% Q1
%%%%%%%%%%%%%%%%%%%%%%%%%%%%%%%%%%%%%%%%%%%%

\textbf{Q1:} Determine the unknown particle:
\[
\bar{d} \rightarrow \bar{u} + x
\]

\vspace{0.5cm}

\begin{tabular}{l l}

% LEFT COLUMN
\begin{minipage}[t]{0.46\textwidth}
\vspace{0pt}
\textbf{Feynman diagram}

\vspace{0.6cm}

% --- TEMPLATE (labels changed only) ---
\begin{tikzpicture}
\begin{feynman}
  % Vertices
  \vertex (l1) at (0,0) {\(\bar{d}\)};
  \vertex [coordinate] (a) at (1,2) {};
  \vertex (l2) at (0,4) {\(\bar{u}\)};
  \vertex (b) at (4,2) {};

  \diagram*{
    (l1) -- [fermion] (a) -- [fermion] (l2),
    (a) -- [boson, edge label=\(x\), momentum'={}] (b),
  };
\end{feynman}
\end{tikzpicture}
% --- END TEMPLATE ---
\end{minipage}
&
% RIGHT COLUMN
\begin{minipage}[t]{0.50\textwidth}
\vspace{0pt}

\renewcommand{\arraystretch}{1.4}
\begin{tabular}{@{}l c c c c c@{}}
 & \(\bar{d}\) & \(\Rightarrow\) & \(\bar{u}\) & \(+\) & \(x\) \\
Charge &  &  &  &  &  \\
Lepton no. &  &  &  &  &  \\
Baryon no. &  &  &  &  &  \\
\end{tabular}

\end{minipage}

\end{tabular}

\vspace{0.5cm}

\textbf{Conclusion:}

\vspace{2.2cm}

%%%%%%%%%%%%%%%%%%%%%%%%%%%%%%%%%%%%%%%%%%%%
% Q2
%%%%%%%%%%%%%%%%%%%%%%%%%%%%%%%%%%%%%%%%%%%%
\newpage
\textbf{Q2:} Determine the unknown particle:
\[
c \rightarrow x + W^+
\]

\vspace{0.5cm}

\begin{tabular}{l l}

% LEFT COLUMN
\begin{minipage}[t]{0.46\textwidth}
\vspace{0pt}
\textbf{Feynman diagram}

\vspace{0.6cm}

% --- TEMPLATE (labels changed only) ---
\begin{tikzpicture}
\begin{feynman}
  % Vertices
  \vertex (l1) at (0,0) {\(c\)};
  \vertex [coordinate] (a) at (1,2) {};
  \vertex (l2) at (0,4) {\(x\)};
  \vertex (b) at (4,2) {};

  \diagram*{
    (l1) -- [fermion] (a) -- [fermion] (l2),
    (a) -- [boson, edge label=\(W^+\), momentum'={}] (b),
  };
\end{feynman}
\end{tikzpicture}
% --- END TEMPLATE ---
\end{minipage}
&
% RIGHT COLUMN
\begin{minipage}[t]{0.50\textwidth}
\vspace{0pt}

\renewcommand{\arraystretch}{1.4}
\begin{tabular}{@{}l c c c c c@{}}
 & \(c\) & \(\Rightarrow\) & \(x\) & \(+\) & \(W^+\) \\
Charge &  &  &  &  &  \\
Lepton no. &  &  &  &  &  \\
Baryon no. &  &  &  &  &  \\
\end{tabular}

\end{minipage}

\end{tabular}

\vspace{0.5cm}

\textbf{Conclusion:}

\vspace{2.2cm}

%%%%%%%%%%%%%%%%%%%%%%%%%%%%%%%%%%%%%%%%%%%%
% Q3
%%%%%%%%%%%%%%%%%%%%%%%%%%%%%%%%%%%%%%%%%%%%

\textbf{Q3:} Determine the unknown particles:
\[
\mu^- \rightarrow x + y
\]

\vspace{0.5cm}

\begin{tabular}{l l}

% LEFT COLUMN
\begin{minipage}[t]{0.46\textwidth}
\vspace{0pt}
\textbf{Feynman diagram}

\vspace{0.6cm}

% --- TEMPLATE (labels changed only) ---
\begin{tikzpicture}
\begin{feynman}
  % Vertices
  \vertex (l1) at (0,0) {\(\mu^-\)};
  \vertex [coordinate] (a) at (1,2) {};
  \vertex (l2) at (0,4) {\(x\)};
  \vertex (b) at (4,2) {};

  \diagram*{
    (l1) -- [fermion] (a) -- [fermion] (l2),
    (a) -- [boson, edge label=\(y\), momentum'={}] (b),
  };
\end{feynman}
\end{tikzpicture}
% --- END TEMPLATE ---
\end{minipage}
&
% RIGHT COLUMN
\begin{minipage}[t]{0.50\textwidth}
\vspace{0pt}

\renewcommand{\arraystretch}{1.4}
\begin{tabular}{@{}l c c c c c@{}}
 & \(\mu^-\) & \(\Rightarrow\) & \(x\) & \(+\) & \(y\) \\
Charge &  &  &  &  &  \\
Lepton no. &  &  &  &  &  \\
Baryon no. &  &  &  &  &  \\
\end{tabular}

\end{minipage}

\end{tabular}

\vspace{0.5cm}

\textbf{Conclusion:}

\vspace{2.2cm}

%%%%%%%%%%%%%%%%%%%%%%%%%%%%%%%%%%%%%%%%%%%%
% Q4
%%%%%%%%%%%%%%%%%%%%%%%%%%%%%%%%%%%%%%%%%%%%
\newpage
\textbf{Q4:} Draw a Feynman diagram to show the initial stage of a valid anti-bottom decay.

\vspace{0.4cm}
\[
\bar{b} \rightarrow \;\;\; x \;+\; y
\]

\vspace{0.5cm}

\begin{tabular}{l l}

% LEFT COLUMN (blank drawing space)
\begin{minipage}[t]{0.46\textwidth}
\vspace{0pt}
\textbf{Feynman diagram}

\vspace{0.6cm}

\begin{tikzpicture}
  \path (0,0) rectangle (6,4.5);
\end{tikzpicture}

\end{minipage}
&
% RIGHT COLUMN (conservation table)
\begin{minipage}[t]{0.50\textwidth}
\vspace{0pt}

\renewcommand{\arraystretch}{1.4}
\begin{tabular}{@{}l c c c c c@{}}
 & \(\bar{b}\) & \(\Rightarrow\) & \(x\) & \(+\) & \(y\) \\
Charge &  &  &  &  &  \\
Lepton no. &  &  &  &  &  \\
Baryon no. &  &  &  &  &  \\
\end{tabular}

\end{minipage}

\end{tabular}

\vspace{0.5cm}

\textbf{Conclusion:}

\vspace{2.2cm}

%%%%%%%%%%%%%%%%%%%%%%%%%%%%%%%%%%%%%%%%%%%%
% Q5
%%%%%%%%%%%%%%%%%%%%%%%%%%%%%%%%%%%%%%%%%%%%

\textbf{Q5:} Draw a Feynman diagram to show the initial stage of a valid tau decay.

\vspace{0.4cm}
\[
\tau^- \rightarrow \;\;\; x \;+\; y
\]

\vspace{0.5cm}

\begin{tabular}{l l}

% LEFT COLUMN (blank drawing space)
\begin{minipage}[t]{0.46\textwidth}
\vspace{0pt}
\textbf{Feynman diagram}

\vspace{0.6cm}

\begin{tikzpicture}
  \path (0,0) rectangle (6,4.5);
\end{tikzpicture}

\end{minipage}
&
% RIGHT COLUMN (conservation table)
\begin{minipage}[t]{0.50\textwidth}
\vspace{0pt}

\renewcommand{\arraystretch}{1.4}
\begin{tabular}{@{}l c c c c c@{}}
 & \(\tau^-\) & \(\Rightarrow\) & \(x\) & \(+\) & \(y\) \\
Charge &  &  &  &  &  \\
Lepton no. &  &  &  &  &  \\
Baryon no. &  &  &  &  &  \\
\end{tabular}

\end{minipage}

\end{tabular}

\vspace{0.5cm}

\textbf{Conclusion:}

\vspace{2.2cm}

%%%%%%%%%%%%%%%%%%%%%%%%%%%%%%%%%%%%%%%%%%%%
% Q6
%%%%%%%%%%%%%%%%%%%%%%%%%%%%%%%%%%%%%%%%%%%%
\newpage
\textbf{Q6:} Draw a Feynman diagram to show the initial stage of a valid top decay.

\vspace{0.4cm}
\[
t \rightarrow \;\;\; x \;+\; y
\]

\vspace{0.5cm}

\begin{tabular}{l l}

% LEFT COLUMN (blank drawing space)
\begin{minipage}[t]{0.46\textwidth}
\vspace{0pt}
\textbf{Feynman diagram}

\vspace{0.6cm}

\begin{tikzpicture}
  \path (0,0) rectangle (6,4.5);
\end{tikzpicture}


\end{minipage}
&
% RIGHT COLUMN (conservation table)
\begin{minipage}[t]{0.50\textwidth}
\vspace{0pt}

\renewcommand{\arraystretch}{1.4}
\begin{tabular}{@{}l c c c c c@{}}
 & \(t\) & \(\Rightarrow\) & \(x\) & \(+\) & \(y\) \\
Charge &  &  &  &  &  \\
Lepton no. &  &  &  &  &  \\
Baryon no. &  &  &  &  &  \\
\end{tabular}

\end{minipage}

\end{tabular}

\vspace{0.5cm}

\textbf{Conclusion:}

\vspace{2.2cm}

\newpage

\section*{Weak Interaction -- Initial Stage \; (Answer Key)}

\textbf{Note:} For Q4--Q6 there are multiple valid answers. One model answer is shown for each, followed by other acceptable options.

\vspace{0.7cm}

%%%%%%%%%%%%%%%%%%%%%%%%%%%%%%%%%%%%%%%%%%%%
% Q1
%%%%%%%%%%%%%%%%%%%%%%%%%%%%%%%%%%%%%%%%%%%%
\textbf{Q1:} \(\bar{d} \rightarrow \bar{u} + x\)

\[
\boxed{x = W^+}
\]

\vspace{0.4cm}

\begin{tabular}{l l}

% LEFT COLUMN
\begin{minipage}[t]{0.46\textwidth}
\vspace{0pt}
\textbf{Feynman diagram}

\vspace{0.6cm}

% --- TEMPLATE (labels changed only) ---
\begin{tikzpicture}
\begin{feynman}
  % Vertices
  \vertex (l1) at (0,0) {\(\bar{d}\)};
  \vertex [coordinate] (a) at (1,2) {};
  \vertex (l2) at (0,4) {\(\bar{u}\)};
  \vertex (b) at (4,2) {};

  \diagram*{
    (l1) -- [fermion] (a) -- [fermion] (l2),
    (a) -- [boson, edge label=\(W^+\), momentum'={}] (b),
  };
\end{feynman}
\end{tikzpicture}
% --- END TEMPLATE ---
\end{minipage}
&
% RIGHT COLUMN
\begin{minipage}[t]{0.50\textwidth}
\vspace{0pt}

\renewcommand{\arraystretch}{1.4}
\begin{tabular}{@{}l c c c c c@{}}
 & \(\bar{d}\) & \(\Rightarrow\) & \(\bar{u}\) & \(+\) & \(W^+\) \\
Charge & \(+\tfrac{1}{3}\) & \(=\) & \(-\tfrac{2}{3}\) & \(+\) & \(+1\) \\
Lepton no. & \(0\) & \(=\) & \(0\) & \(+\) & \(0\) \\
Baryon no. & \(-\tfrac{1}{3}\) & \(=\) & \(-\tfrac{1}{3}\) & \(+\) & \(0\) \\
\end{tabular}

\end{minipage}

\end{tabular}

\vspace{0.8cm}

%%%%%%%%%%%%%%%%%%%%%%%%%%%%%%%%%%%%%%%%%%%%
% Q2
%%%%%%%%%%%%%%%%%%%%%%%%%%%%%%%%%%%%%%%%%%%%
\newpage
\textbf{Q2:} \(c \rightarrow x + W^+\)

\[
\boxed{x = s \;\; \text{(also valid: } x=d\text{)}}
\]

\vspace{0.4cm}

\begin{tabular}{l l}

% LEFT COLUMN
\begin{minipage}[t]{0.46\textwidth}
\vspace{0pt}
\textbf{Feynman diagram}

\vspace{0.6cm}

% --- TEMPLATE (labels changed only) ---
\begin{tikzpicture}
\begin{feynman}
  % Vertices
  \vertex (l1) at (0,0) {\(c\)};
  \vertex [coordinate] (a) at (1,2) {};
  \vertex (l2) at (0,4) {\(s\)};
  \vertex (b) at (4,2) {};

  \diagram*{
    (l1) -- [fermion] (a) -- [fermion] (l2),
    (a) -- [boson, edge label=\(W^+\), momentum'={}] (b),
  };
\end{feynman}
\end{tikzpicture}
% --- END TEMPLATE ---
\end{minipage}
&
% RIGHT COLUMN
\begin{minipage}[t]{0.50\textwidth}
\vspace{0pt}

\renewcommand{\arraystretch}{1.4}
\begin{tabular}{@{}l c c c c c@{}}
 & \(c\) & \(\Rightarrow\) & \(s\) & \(+\) & \(W^+\) \\
Charge & \(+\tfrac{2}{3}\) & \(=\) & \(-\tfrac{1}{3}\) & \(+\) & \(+1\) \\
Lepton no. & \(0\) & \(=\) & \(0\) & \(+\) & \(0\) \\
Baryon no. & \(+\tfrac{1}{3}\) & \(=\) & \(+\tfrac{1}{3}\) & \(+\) & \(0\) \\
\end{tabular}

\end{minipage}

\end{tabular}

\vspace{0.8cm}

%%%%%%%%%%%%%%%%%%%%%%%%%%%%%%%%%%%%%%%%%%%%
% Q3
%%%%%%%%%%%%%%%%%%%%%%%%%%%%%%%%%%%%%%%%%%%%
\textbf{Q3:} \(\mu^- \rightarrow x + y\)

\[
\boxed{x=\nu_\mu \quad\text{and}\quad y=W^-}
\]

\vspace{0.4cm}

\begin{tabular}{l l}

% LEFT COLUMN
\begin{minipage}[t]{0.46\textwidth}
\vspace{0pt}
\textbf{Feynman diagram}

\vspace{0.6cm}

% --- TEMPLATE (labels changed only) ---
\begin{tikzpicture}
\begin{feynman}
  % Vertices
  \vertex (l1) at (0,0) {\(\mu^-\)};
  \vertex [coordinate] (a) at (1,2) {};
  \vertex (l2) at (0,4) {\(\nu_\mu\)};
  \vertex (b) at (4,2) {};

  \diagram*{
    (l1) -- [fermion] (a) -- [fermion] (l2),
    (a) -- [boson, edge label=\(W^-\), momentum'={}] (b),
  };
\end{feynman}
\end{tikzpicture}
% --- END TEMPLATE ---
\end{minipage}
&
% RIGHT COLUMN
\begin{minipage}[t]{0.50\textwidth}
\vspace{0pt}

\renewcommand{\arraystretch}{1.4}
\begin{tabular}{@{}l c c c c c@{}}
 & \(\mu^-\) & \(\Rightarrow\) & \(\nu_\mu\) & \(+\) & \(W^-\) \\
Charge & \(-1\) & \(=\) & \(0\) & \(+\) & \(-1\) \\
Lepton no. & \(+1\) & \(=\) & \(+1\) & \(+\) & \(0\) \\
Baryon no. & \(0\) & \(=\) & \(0\) & \(+\) & \(0\) \\
\end{tabular}

\end{minipage}

\end{tabular}

\vspace{1.0cm}

%%%%%%%%%%%%%%%%%%%%%%%%%%%%%%%%%%%%%%%%%%%%
% Q4
%%%%%%%%%%%%%%%%%%%%%%%%%%%%%%%%%%%%%%%%%%%%
\newpage
\textbf{Q4:} Draw a Feynman diagram to show the initial stage of a valid anti-bottom decay.

\vspace{0.3cm}
\textbf{Model valid answer:}
\[
\boxed{\bar{b} \rightarrow \bar{c} + W^+}
\]

\vspace{0.4cm}

\begin{tikzpicture}
\begin{feynman}
  \vertex (l1) at (0,0) {\(\bar{b}\)};
  \vertex [coordinate] (a) at (1,2) {};
  \vertex (l2) at (0,4) {\(\bar{c}\)};
  \vertex (b) at (4,2) {};

  \diagram*{
    (l1) -- [fermion] (a) -- [fermion] (l2),
    (a) -- [boson, edge label=\(W^+\)] (b),
  };
\end{feynman}
\end{tikzpicture}

\vspace{0.5cm}
\textbf{Other acceptable answers include:}
\[
\bar{b}\rightarrow \bar{u}+W^+ \quad (\text{due to energy considerations: } \bar{b}\rightarrow \bar{t}+W^+ \text{ is not allowed})
\]

\vspace{0.8cm}

%%%%%%%%%%%%%%%%%%%%%%%%%%%%%%%%%%%%%%%%%%%%
% Q5
%%%%%%%%%%%%%%%%%%%%%%%%%%%%%%%%%%%%%%%%%%%%
\textbf{Q5:} Draw a Feynman diagram to show the initial stage of a valid tau decay.

\vspace{0.3cm}
\textbf{Model valid answer:}
\[
\boxed{\tau^- \rightarrow \nu_\tau + W^-}
\]

\vspace{0.4cm}

\begin{tikzpicture}
\begin{feynman}
  \vertex (l1) at (0,0) {\(\tau^-\)};
  \vertex [coordinate] (a) at (1,2) {};
  \vertex (l2) at (0,4) {\(\nu_\tau\)};
  \vertex (b) at (4,2) {};

  \diagram*{
    (l1) -- [fermion] (a) -- [fermion] (l2),
    (a) -- [boson, edge label=\(W^-\)] (b),
  };
\end{feynman}
\end{tikzpicture}

\vspace{0.5cm}
\textbf{No other acceptable answers in this case}

\vspace{0.8cm}

%%%%%%%%%%%%%%%%%%%%%%%%%%%%%%%%%%%%%%%%%%%%
% Q6
%%%%%%%%%%%%%%%%%%%%%%%%%%%%%%%%%%%%%%%%%%%%
\newpage
\textbf{Q6:} Draw a Feynman diagram to show the initial stage of a valid top decay.

\vspace{0.3cm}
\textbf{Model valid answer:}
\[
\boxed{t \rightarrow b + W^+}
\]

\vspace{0.4cm}

\begin{tikzpicture}
\begin{feynman}
  \vertex (l1) at (0,0) {\(t\)};
  \vertex [coordinate] (a) at (1,2) {};
  \vertex (l2) at (0,4) {\(b\)};
  \vertex (b) at (4,2) {};

  \diagram*{
    (l1) -- [fermion] (a) -- [fermion] (l2),
    (a) -- [boson, edge label=\(W^+\)] (b),
  };
\end{feynman}
\end{tikzpicture}

\vspace{0.5cm}
\textbf{Other acceptable answers include:}
\[
t\rightarrow s+W^+ \quad \text{or} \quad t\rightarrow d+W^+
\]

\end{document}
